\section{Distribution and scaling}
\label{discussion:distribution-and-scaling}

This section reads the distributed results of Section~\ref{results:distributed-join-and-scaling} against the distribution mechanisms of Chapter~\ref{ch:background} to answer the scaling comparison \acrshort{rq}2 poses: why the broadcast join scaled as it did, when distribution earns the coordination cost it imposes, and where a tuned single node remains the better choice. Those results reported how the national-scale spatial join scaled, how it stood against the single-node baselines, how its execution decomposed into phases, and what it cost, without attributing any of those orderings to a cause.

The strongest case for distribution is that on the largest workload it was necessary at all. Neither single-node engine completed the large-tier join, so no single-node baseline is drawn at that tier and only the distributed curve appears in Figure~\ref{fig:distributed-wall-clock-vs-workers}. This is best read as the single-node memory ceiling: a join whose working set exceeds the memory of one machine cannot complete on that machine, and because DuckDB scales only vertically, enlarging one node rather than adding more, the shortfall cannot be absorbed by attaching workers \parencite{duckdb2025_scalability_faq} (Section~\ref{background:cloud-native-data-architecture:distributed-and-non-distributed-computation}). The distinction at the large tier is therefore one of feasibility rather than speed: the single-node engines returned no result while the distributed broadcast join did.

Where both paths completed, the broadcast join scaled better as the work grew. Speedup rose monotonically in worker count at every tier, yet stayed below the ideal-linear reference and fell further short as workers were added (Figure~\ref{fig:distributed-speedup-efficiency}). Distribution repays its coordination cost only once the work is large enough to amortize it \parencite{mcsherry2015_scalability_but_at_what_cost,zaharia2012_resilient_distributed_datasets}, and the widening shortfall from ideal points to a fixed, worker-independent overhead. Driver-side planning, together with the four coordination sources of task scheduling, data shuffling, serialization, and network communication, does not shrink as executors are added, so each added worker returns less marginal speedup \parencite{mcsherry2015_scalability_but_at_what_cost} (Section~\ref{background:cloud-native-data-architecture:distributed-and-non-distributed-computation}). The paired parallel efficiency reads the same relation from the other side: it is lowest where that fixed cost is largest relative to the work, collapsing at the small tier and holding high at the large tier.

The phase breakdown locates where that limit lies. Executor run time accounted for effectively the whole of the stacked phase time at every tier and worker count, while driver collection contributed a negligible sliver, recorded as zero at several worker counts (Figure~\ref{fig:distributed-execution-phase-breakdown}). The binding work is therefore executor-side, not result collection at the driver, and that executor phase did not fall as workers were added; the shuffle volume behaved likewise, growing with tier but varying only modestly with worker count within a tier. A job limited by result collection would have shown the opposite split, with the driver rather than the executors dominating. The pattern instead matches spatial query skew, in which the realized parallelism is bounded above by the most loaded partition, so once many records or query windows concentrate on one executor that executor becomes a straggler and adding workers cannot push run-time below the time needed to drain the heaviest partition \parencite{tang2020_locationspark_distributed_spatial_query,kwon2012_skewtune_mitigating_skew} (Section~\ref{background:distributed-spatial-processing:spatial-partitioning-and-parallel-efficiency}). That the dominant phase is executor-side and insensitive to worker count points to partition load, rather than driver coordination, as the ceiling on speedup. This is consistent with an efficiency that improves with tier yet settles below one even where the work is large.

Network transfer on the distributed path carries a meaning of its own, and reading it requires separating two quantities the Results chapter keeps apart. At the Sedona driver boundary the bytes received exceeded the bytes sent by roughly a factor of two at every tier, and both directions grew only modestly with tier, remaining at a few tens of kilobytes (Figure~\ref{fig:distributed-client-boundary}). This is the broadcast-and-collect fingerprint: the small input is broadcast out to the executors and a small result set is collected back, so the volume crossing the driver boundary is governed by the size of the result and the broadcast input, not by the size of the data the join scans \parencite{yu2019_geospark_perspective_beyond} (Section~\ref{background:distributed-spatial-processing:the-lifecycle-of-a-distributed-spatial-query}). This boundary transfer must not be conflated with the shuffle volume of the phase breakdown, the executor-to-executor redistribution measured in gigabytes and larger by orders of magnitude (Figure~\ref{fig:distributed-execution-phase-breakdown}); the two are different channels read at different points. The driver-boundary transfer is the only transfer on this path that reaches operational cost, through the cross-region egress term of the cost model (Equation~\ref{eq:cost-model}, Section~\ref{research-design-and-methodology:metrics-and-cost-model:cost-model}). The distributed configuration is the only one to carry a non-zero egress term, yet the volume that flows through it is small, so its contribution to cost is small in turn. The byte signature here records where the result is assembled and returned. Like the host-boundary count throughout, it is read per architecture, so it is not a transfer ranking against the single-node configurations, whose own host-boundary reading is established in Section~\ref{discussion:single-node-cloud-native-vs-traditional}.

Against the single-node baselines, the worker count at which the broadcast curve overtakes a baseline is where distribution begins to pay (Section~\ref{research-design-and-methodology:statistical-analysis:answering-rq2-scaling-against-single-node}). At the small tier it never does within the tested range: the broadcast join fell with added workers and approached the PostGIS minimum without dropping below it across the tested \numrange{2}{16} workers (Figure~\ref{fig:distributed-wall-clock-vs-workers}). At \num{16} workers the broadcast confidence interval still overlaps the PostGIS minimum, and the Results chapter declares no winner there. Any worker count at which the small-tier curve would cross the PostGIS baseline lies beyond the \numrange{2}{16} workers tested, so that crossover is an extrapolation rather than a measured point. Within the tested range, a tuned single node remains the better choice for a small join. The medium tier behaves in reverse: the broadcast join was already far below the fastest single-node engine at \num{2} workers and led at every worker count tested, so for work of that size distribution pays from the smallest cluster measured. The contrast across tiers, rather than any single crossover, is what locates the size at which distribution becomes worthwhile.

The choice of join strategy was the largest single lever on this path. At the small tier, the only tier at which it completed, the partitioned join ran roughly two orders of magnitude slower than the broadcast join and cost roughly \num{30} times as much; at the medium and large tiers it did not complete at any worker count, recording an executor out-of-memory failure throughout (Figure~\ref{fig:distributed-strategy-contrast}, Table~\ref{tab:distributed-join-descriptive-statistics}). The penalty follows from what each strategy moves. A partitioned join repartitions both inputs through a shuffle so that spatially matching records co-locate, and because a polygon may straddle partition boundaries it is duplicated into every overlapping partition to preserve correctness, which inflates both the shuffle volume and the per-executor working set \parencite{yu2019_geospark_perspective_beyond} (Section~\ref{background:distributed-spatial-processing:distributed-spatial-operations}). A broadcast join avoids the key-based redistribution, shipping only the small input to each executor, so its strict shuffle cost stays low when one input is small. At the sizes tested here one input was always small enough to broadcast, so the redistribution the partitioned strategy paid for bought no offsetting benefit, and the duplicated boundary geometries are the most likely source of the executor memory exhaustion that ended the medium and large runs.